\section{Discussion}
\label{sec:discussion}

This section gathers the interpretations the results support, each at the strength its
evidence carries, and names the measurement that would settle it where one is known.

\paragraph{One mechanism, indexed per stratum.} Every measurable stratum carries a diagonal
clock on its own local group (Section~\ref{sec:strata-result}). Strata that share a local
group select exactly the same characters of it in $158$ of $160$ pre-registered pairs
(Section~\ref{sec:strata-sameness}). At $n = 113$ and $n = 121$, deleting the key characters
from the embedding returns the network to chance (Section~\ref{sec:causal-internal}). The
reading these support is that the network runs the same character-based clock on each
stratum, on that stratum's own local group, and not a private sub-circuit per stratum. The
sameness claim ships without a failed-run control, because the project has only two scorable
failed-run pairs and both come from one model. The same statistic on enough models that never
generalised would complete it.

\paragraph{What the nilpotents do.} \citet{chen2026multiplication} describe non-regular
$\mathcal{J}$-classes as breaking the reduction to local group characters. At every
non-square-free modulus we measured, the reduction does not break. A stratum collapses onto a
smaller local group, $G_{m}$ of Theorem~\ref{thm:stratum}, and the clock runs there
(Sections~\ref{sec:clock-primepowers} and \ref{sec:strata-result}). The coordinate that makes
this measurable was first sighted at one seed (1 seed), and the five-seed measurement of
Section~\ref{sec:strata} supersedes that sighting. The ablation reaches nilpotency depth $7$
at $n = 2^{7}$, and depth beyond that is untested (Section~\ref{sec:causal-128}).

\paragraph{Where the additive structure comes from.} The key additive frequencies of a
grokked model are exactly $P(n)$, one family per maximal prime power
(Section~\ref{sec:crt-law}). The $\mathcal{J}$-class indicators alone concentrate their
energy on the duals of every divisor, and select $P(n)$ at only $2$ of the $10$ composite
moduli tested. The selection of the maximal-prime-power subfamily
is therefore the trained network's, since the algebra alone does not make it. Why the network
makes that selection is not addressed here.

\paragraph{Relation to other mechanistic accounts.} On the unit stratum of a prime the
mechanism is \citet{nanda2023progress}'s clock in discrete-log coordinates, and our engine
reproduces their addition result (Section~\ref{sec:calib}). For $a \cdot b \bmod n$ the two
levels of universality that \citet{chughtai2023toy} describe are algebraic: the family is the
character clock, and the index set is fixed by the stratum's own local group. The top logit
components are $a+b$ terms in discrete-log coordinates at five seeds
(Section~\ref{sec:clock-circuit}), at the one hyperparameter setting we train.
\citet{zhong2023clock} show that other settings can yield other algorithms, and we do not
test them. The basis dependence of \citet{nguyen2026dlogclock} replicates, and the agreement
is specific to float32 (Section~\ref{sec:calib}). The logit ablation survives the three
initialisation conditions of \citet{kunin2024getrich} (Section~\ref{sec:causal-init}).
Section~\ref{sec:limits} names the next test, the same measurement on non-commutative finite
groups.

\paragraph{Precision and absolute numbers.} Training precision changes the measured sparsity
in PyTorch: the same code reads $0.5791$ at float32 and $0.7728$ at float64
(Section~\ref{sec:impl}). A published Gini is therefore comparable to another only at a named
precision, and every absolute sparsity number this paper reports from PyTorch is a float32 measurement.
Comparisons between moduli, and comparisons within one arm, are unaffected. Why our engine
reaches the float64 answer at either dtype is open, with no surviving candidate, and an
operator-by-operator differential trace of one training step at both dtypes would settle it.

\paragraph{What the negatives say about progress measures.} Two statistics that look like
progress measures fire on runs that never generalised. G1, the criterion our pre-registration
named primary, passes on $20$ of $20$ failed stratum-measurements
(Section~\ref{sec:strata-falsifier}). The early CRT enrichment ramp rises monotonically in
failed runs and ends above the band grokked runs occupy early (Section~\ref{sec:crt-ramp}). A
statistic that fires on failures is not evidence, whatever it does on successes. The
procedure we draw from both is to score a candidate statistic on a failed-run control before
proposing it.

\paragraph{Why the shape of this result may travel further than the task.} What the
measurements describe is a network that partitions its input space along a structure the
\emph{data} has, runs one mechanism on every part, and indexes that mechanism by each
part's own symmetry group. Nothing in that sentence mentions modular arithmetic. It is the
pattern an architecture exhibits whenever it routes inputs to sub-computations, as a mixture
of experts does when it chooses among them, or a network whose layers are built to respect a
group action. In those settings the partition is imposed by the designer, where here it was
recovered from a trained model that was never told the algebra of $n$. The instrument transfers with
it: a stratification anyone can compute from the task, a per-stratum statistic, a criterion
fixed before the numbers exist, and a control that has the stratification's size structure
without its content. That last element is what killed three of our own claims
(Appendix~\ref{app:retractions}). We make no claim that a larger model does this; we claim
that this is a testable thing to ask of one.
